\documentclass[UTF8,fontset=fandol,11pt]{ctexart}
\usepackage[a4paper,margin=25mm]{geometry}
\usepackage{amsmath,amssymb,amsthm,mathtools}
\usepackage{xcolor}
\usepackage{booktabs,array,enumitem,fancyhdr,hyperref}
\hypersetup{colorlinks=true,linkcolor=blue!50!black,urlcolor=blue!50!black}
\setlength{\emergencystretch}{2em}
\setlength{\parskip}{0.3em}
\linespread{1.16}
\allowdisplaybreaks
\newtheorem{theorem}{定理}[section]
\newtheorem{proposition}[theorem]{命题}
\newtheorem{lemma}[theorem]{引理}
\theoremstyle{definition}
\newtheorem{definition}[theorem]{定义}
\theoremstyle{remark}
\newtheorem{remark}[theorem]{说明}
\newcommand{\dd}{\,\mathrm d}
\newcommand{\ii}{\mathrm i}
\newcommand{\Rea}{\operatorname{Re}}
\newcommand{\Ima}{\operatorname{Im}}
\newcommand{\Li}{\operatorname{Li}}
\newcommand{\supp}{\operatorname{supp}}
\pagestyle{fancy}
\fancyhf{}
\fancyhead[L]{\small 学习笔记}
\fancyhead[R]{\small 初稿}
\fancyfoot[C]{\thepage}
\setlength{\headheight}{15pt}
\author{}
\date{2026年10月22日初稿}
\makeatletter
\renewcommand{\maketitle}{%
  \thispagestyle{plain}%
  \begin{center}
    {\LARGE\bfseries \@title\par}
    \vspace{0.6em}
    {\small \@date\par}
  \end{center}
  \vspace{0.6em}}
\makeatother
\newcommand{\draftnotice}{\noindent{\small
本文为供逐篇审阅的学习笔记初稿，不是新增研究成果。
文中区分直接证明与引用结果；正文不依赖本文。
本轮已作公式、条件及排版初步核查，尚待作者逐句审定。}\par\medskip}

\usepackage{pgfplots}
\usepackage{float}
\usetikzlibrary{arrows.meta}
\pgfplotsset{compat=1.18}
\date{2026年10月22日初稿；10月24日增补示意图}
\title{振荡相消能带来多大改善}
\begin{document}
\maketitle
\draftnotice
\section{为什么逐项取绝对值会损失信息}
若 $N$ 个复数各有模 $1$，三角不等式给出和的模不超过 $N$。
这个界总是成立，但它忽略了方向。方向规律可以使实际和有界，
也可以在某些参数下使所有项几乎同向。理解振荡相消，首先要确定相消依赖什么条件。

\section{有限几何和的精确计算}
对 $e^{\ii\theta}\ne1$，
\[
 S_N(\theta)=\sum_{k=0}^{N-1}e^{\ii k\theta}
 =\frac{1-e^{\ii N\theta}}{1-e^{\ii\theta}}
 =e^{\ii(N-1)\theta/2}
   \frac{\sin(N\theta/2)}{\sin(\theta/2)}.
\]
第一步来自把和乘以 $1-e^{\ii\theta}$ 后逐项相消；
第二步使用 $1-e^{\ii v}=-2\ii e^{\ii v/2}\sin(v/2)$。
所以
\[
 |S_N(\theta)|\le
 \min\left(N,\frac1{|\sin(\theta/2)|}\right).
\]
若 $\theta$ 固定且不为 $2\pi$ 的整数倍，上界独立于 $N$。
例如 $\theta=\pi$ 时正负交替，和只有 $0$ 或 $1$，
而逐项绝对值的和仍为 $N$。

\section{不能忽略靠近共振的区域}
若 $\theta=0$，所有项同向，$S_N(0)=N$。
更一般地取 $\theta=\alpha/N$，固定 $\alpha$，则当 $\alpha\ne0$ 时
\[
 \frac{|S_N(\alpha/N)|}{N}
 \longrightarrow\left|\frac{2\sin(\alpha/2)}{\alpha}\right|.
\]
在 $\alpha=0$ 处按连续极限取 $1$。
因此即使对每个固定的非共振参数 $\theta$ 都能得到有界结果，也不能把它说成对所有 $\theta$ 一致有界。
参数随 $N$ 靠近共振点时，和仍可能具有 $N$ 的量级。

图~\ref{fig:cancellation-resonance} 将这两个情形放在一起比较。
左图中六个单位向量首尾相接，恰好闭合；右图仍取同一个固定角 $\theta=\pi/3$，
并与随 $N$ 改变的 $\theta=\pi/N$ 对照。
前者的 $|S_N|/N$ 趋于 $0$，后者则趋于 $2/\pi$。
这一区别来自参数是否随项数接近共振点，而不是三角不等式失效。

\begin{figure}[H]
\centering
\begin{minipage}[t]{0.47\linewidth}
\vspace{0pt}
\centering
\begin{tikzpicture}[x=1.55cm,y=1.55cm,>=Stealth,font=\small]
\path[use as bounding box] (-0.85,-0.85) rectangle (2.3,2.65);
\node at (0.7,2.4) {六个方向的精确相消};
\draw[->,gray!70] (-0.8,0)--(2.1,0) node[below] {$\Rea$};
\draw[->,gray!70] (0,-0.4)--(0,2.15) node[left] {$\Ima$};
\coordinate (P0) at (0,0);
\coordinate (P1) at (1,0);
\coordinate (P2) at (1.5,{sqrt(3)/2});
\coordinate (P3) at (1,{sqrt(3)});
\coordinate (P4) at (0,{sqrt(3)});
\coordinate (P5) at (-0.5,{sqrt(3)/2});
\draw[->,very thick,blue!65!black] (P0)--(P1) node[midway,below] {$k=0$};
\draw[->,very thick,blue!65!black] (P1)--(P2) node[midway,right] {$1$};
\draw[->,very thick,blue!65!black] (P2)--(P3) node[midway,right] {$2$};
\draw[->,very thick,blue!65!black] (P3)--(P4) node[midway,above] {$3$};
\draw[->,very thick,blue!65!black] (P4)--(P5) node[midway,left] {$4$};
\draw[->,very thick,blue!65!black] (P5)--(P0) node[midway,left] {$5$};
\fill (P0) circle (1.5pt);
\node[below left] at (P0) {$0$};
\node at (0.65,-0.62) {$S_6(\pi/3)=0$};
\end{tikzpicture}
\end{minipage}\hfill
\begin{minipage}[t]{0.51\linewidth}
\vspace{0pt}
\centering
\begin{tikzpicture}
\begin{axis}[
  width=7.8cm,height=6.3cm,
  xmin=1,xmax=80,ymin=-0.02,ymax=1.08,
  xlabel={$N$},ylabel={$|S_N(\theta)|/N$},
  title={固定参数与接近共振},
  title style={font=\small},label style={font=\small},
  tick label style={font=\small},
  xtick={1,20,40,60,80},ytick={0,0.25,0.5,0.75,1},
  axis lines=left,grid=major,grid style={gray!15},
  legend style={font=\footnotesize,draw=none,
    at={(0.5,-0.23)},anchor=north,legend columns=1},
  domain=1:80,samples=80]
\addplot[black,thick,mark=o,mark size=1.3pt]
  {abs(sin(deg(pi*x/6)))/(x*sin(deg(pi/6)))};
\addlegendentry{固定 $\theta=\pi/3$}
\addplot[blue!65!black,thick,mark=*,mark size=1pt]
  {1/(x*sin(deg(pi/(2*x))))};
\addlegendentry{$\theta=\pi/N$}
\addplot[gray!80!black,dashed,thick,domain=1:80,samples=2] {2/pi};
\addlegendentry{极限 $2/\pi$}
\end{axis}
\end{tikzpicture}
\end{minipage}

\caption{单位向量的相消与接近共振时的不同尺度。
左图取 $N=6$、$\theta=\pi/3$，箭头依次表示 $e^{\ii k\theta}$（$k=0,\ldots,5$），和为 $0$。
右图取整数 $N=1,\ldots,80$；点表示有限和的精确公式值，连线仅便于观察，水平虚线为 $2/\pi$。
归一化后的有限图形与上述极限一致，但极限仍由公式证明。}
\label{fig:cancellation-resonance}
\end{figure}

\section{加入缓变系数：分部求和}
记 $B_n=\sum_{k=1}^n b_k$。有限求和恒等式为
\[
 \sum_{n=1}^Na_nb_n=a_NB_N+
 \sum_{n=1}^{N-1}(a_n-a_{n+1})B_n.
\]
只需在右侧展开 $B_n$，每个 $b_k$ 的系数会望远镜相消为 $a_k$。
因此若 $|B_n|\le C$，便有
\[
 \left|\sum_{n=1}^Na_nb_n\right|
 \le C\left(|a_N|+\sum_{n=1}^{N-1}|a_n-a_{n+1}|\right).
\]
这不仅适用于单调系数，而是由离散总变差控制。
若 $a_n\ge0$ 单调下降，右侧恰为 $Ca_1$。

取 $b_n=e^{\ii n\theta}$，固定 $\theta\notin2\pi\mathbb Z$，其部分和有统一界。
若 $a_n\downarrow0$，对任意尾区间同样分部求和，得到
$\left|\sum_{n=M}^Na_nb_n\right|\le C_\theta a_M$。
这证明了 Cauchy 条件，即 Dirichlet 判别法在本情形下成立。
若系数没有趋零，或没有所需的变差控制，便不能照搬结论。
这里的常数 $C_\theta$ 依赖于与共振点的距离。
若同时让 $\theta$ 趋近 $2\pi\mathbb Z$，必须重新估计这个常数；
固定参数下的尾和判别不自动给出双参数的一致收敛。

\section{一个条件收敛级数的精确值}
固定 $0<\theta<2\pi$。对 $0<r<1$，绝对收敛的幂级数给出
\[
 \sum_{n=1}^\infty\frac{r^ne^{\ii n\theta}}n
 =-\log(1-re^{\ii\theta}).
\]
它可由几何级数逐项积分获得，取从 $r=0$ 连续延伸的对数分支。
取实部得
\[
 \sum_{n=1}^\infty\frac{r^n\cos(n\theta)}n
 =-\log|1-re^{\ii\theta}|.
\]
为令 $r\uparrow1$，注意系数 $r^n/n$ 单调下降，前节尾和界至多为
$C_\theta/(M+1)$，并且对 $0<r\le1$ 一致。
先固定 $M$ 令 $r\uparrow1$，再令 $M\to\infty$，于是
\[
 \sum_{n=1}^\infty\frac{\cos(n\theta)}n
 =-\log\left(2\sin\frac\theta2\right).
\]
该极限交换有尾和界支持；不是对条件收敛级数随意逐项取极限。
当 $\theta\to0$ 时右侧发散，恰与共振处的调和级数对应。

\section{这对零点振荡说明了什么}
在对数变量 $u=\log x$ 中，显式公式含有 $e^{\ii\gamma u}$ 一类项。
对全部项取绝对值可能丢失相位信息；但零点虚部不等间距，
系数也不是固定的几何序列，所以本篇的闭式不能直接用于零点和。
若要利用相消，必须建立适合那些实际频率的部分和估计，
并保证截断、配对和极限的规则合法。

因此“存在振荡”只是提出相消问题的起点，不是已经证明了足够相消。
本篇提供可完全核查的模型，帮助辨认估计中保留或丢失了哪些信息。

\section*{参考依据与后续核验}
几何和、有限分部求和及所需的极限交换均已在文中说明。
可与本书《零点项的对数尺度振荡》对照阅读。
图~\ref{fig:cancellation-resonance} 的绘图源码随本文提供；它同时展示固定参数与接近共振的参数，
不将本篇的等距频率模型直接当作 $\zeta$ 零点和的相消估计。
\end{document}
